\chapter{Identités et fonctions trigonométriques}
\section*{Exercice 15.1 -- Addition d'angles}
La formule de différence donne $\sin(\pi/4)\cos(\pi/6)-\cos(\pi/4)\sin(\pi/6)=(\sqrt6-\sqrt2)/4$. Le résultat est positif, comme l'exige un angle du premier quadrant.

\section*{Exercice 15.2 -- Angle double}
La formule d'addition donne $\cos(2x)=\cos^2x-\sin^2x=1-2\sin^2x$. L'équation équivaut à $\sin x=0$, soit $x=k\pi$ pour $k\in\Z$.

\section*{Exercice 15.3 -- Dériver un signal}
L'amplitude est $3$ et la période $2\pi/2=\pi$. La composition donne $f'(t)=-6\sin(2t+\pi/6)$ puis $f''(t)=-12\cos(2t+\pi/6)=-4f(t)$. Le facteur $2$ intervient à chaque dérivation.

\section*{Exercice 15.4 -- Rotation}
La formule donne $(x,y)\mapsto(-y,x)$, donc $(2,-3)\mapsto(3,2)$. Les normes au carré valent toutes deux $13$. La conservation provient de l'identité $\cos^2\theta+\sin^2\theta=1$, pas d'une coïncidence numérique.

\section*{Exercice 15.5 -- Produit en somme}
Soustraire les deux formules du cosinus annule $\cos a\cos b$ et laisse $2\sin a\sin b$. Pour $a=b=x$, on obtient $2\sin^2x=1-\cos2x$, utile pour transformer une puissance en oscillation de fréquence double.

\section*{Exercice 15.6 -- Limite fondamentale}
Lorsque $h\to0^+$, $\cos h\to1$ ; les gendarmes donnent $\sin h/h\to1$. La fonction $h\mapsto\sin h/h$ est paire sur son domaine car sinus et dénominateur changent tous deux de signe. La limite à gauche est donc la même. Le radian est essentiel à l'encadrement utilisé.
